% ===============================
% Worksheet / Manual Template
% ===============================
\documentclass[12pt, a4paper]{article}

% ===============================
% Metadata
% ===============================
\newcommand*{\CourseCode}{MAT221}
\newcommand*{\CourseTitle}{Real Analysis}
\newcommand*{\Chapter}{Riemann Integral}
\newcommand*{\Topics}{Upper and Lower Sums, Riemann Integrability}
\newcommand*{\DocDate}{December 30, 2025}

\include{header}

\begin{document}
\FrontPage
\Large

% ======================================================
\begin{heading}
Historical and Conceptual Motivation
\end{heading}

\begin{quote}
\textit{``The integral was not born as a limit—it was born as area.''}
\end{quote}

In elementary calculus, integration is often introduced as the inverse operation of differentiation.
That approach is historically inaccurate and mathematically insufficient.

The founders of calculus—Newton and Leibniz—viewed integration primarily through antiderivatives.
However, this immediately leads to a serious limitation:

\clearpage

\begin{problem}
Show that the function
\[
h(x)=
\begin{cases}
1, & 0\le x<1,\\
2, & 1\le x\le 2,
\end{cases}
\]
has no antiderivative on $[0,2]$.
\end{problem}

\vfill

Despite this, it is intuitively clear that the area under $h$ from $0$ to $2$ should be $3$.
Thus, **area must be defined independently of derivatives**.

This philosophical shift—separating integration from differentiation—was made rigorous by
Cauchy and later refined by Riemann around 1850 :contentReference[oaicite:1]{index=1}.

\clearpage

% ======================================================
\begin{heading}
From Geometry to Approximation
\end{heading}

The geometric idea behind integration is approximation:

\begin{center}
\emph{Approximate the region under a curve by rectangles, and refine the approximation.}
\end{center}

Each approximation is crude, but refinement improves accuracy.
The mathematical challenge is to turn this idea into a precise definition that depends
only on properties of real numbers—especially completeness.

\clearpage

% ======================================================
\begin{heading}
Bounded Functions on Closed Intervals
\end{heading}

\begin{definition}[Bounded Function]
A function $f:[a,b]\to\mathbb{R}$ is called \textbf{bounded} if there exists $M>0$ such that
\[
|f(x)|\le M \quad \text{for all } x\in[a,b].
\]
\end{definition}

Boundedness is not optional.
If $f$ is unbounded on any subinterval, then no finite rectangular approximation is possible.

\vspace{100pt}

\begin{theorem}
Every continuous function on a closed interval $[a,b]$ is bounded.
\end{theorem}

This is a direct consequence of compactness and the Extreme Value Theorem.

\clearpage

% ======================================================
\begin{heading}
Partitions of an Interval
\end{heading}

\begin{definition}[Partition]
A \textbf{partition} $P$ of $[a,b]$ is a finite ordered set
\[
P=\{x_0,x_1,\dots,x_n\}
\]
such that
\[
a=x_0<x_1<\cdots<x_n=b.
\]
\end{definition}

Each partition induces subintervals $[x_{i-1},x_i]$.

\clearpage

\begin{definition}[Norm (Mesh)]
The \textbf{norm} of a partition $P$ is
\[
\|P\|=\max_{1\le i\le n}(x_i-x_{i-1}).
\]
\end{definition}

A small norm corresponds to a fine partition.
In the limit, we want $\|P\|\to 0$.

\clearpage

% ======================================================
\begin{heading}
Supremum and Infimum on Subintervals
\end{heading}

\begin{definition}
Let $f$ be bounded function on $[a,b]$ and let $[x_{i-1},x_i]$ be a subinterval of the interval. Define
\[M_i=\sup\{f(x):x\in[x_{i-1},x_i]\}.\]
\[m_i=\inf\{f(x):x\in[x_{i-1},x_i]\}.\]
\end{definition}

These quantities always exist because:
\begin{itemize}
\item $f$ is bounded,
\item $\mathbb{R}$ is complete.
\end{itemize}

They represent the best possible overestimate and underestimate of $f$ on each subinterval.

\clearpage

% ======================================================
\begin{heading}
Upper and Lower Sums
\end{heading}

\begin{definition}[Lower Sum]
Let $f$ be bounded function on $[a,b]$ and let $P=\{x_0,x_1,\dots,x_n\}$ be a partition of $[a,b]$. Define
\[L(f,P)=\sum_{i=1}^n m_i(x_i-x_{i-1}).\]
\end{definition}

\clearpage

\begin{definition}[Upper Sum]
Let $f$ be bounded function on $[a,b]$ and let $P=\{x_0,x_1,\dots,x_n\}$ be a partition of $[a,b]$. Define
\[U(f,P)=\sum_{i=1}^n M_i(x_i-x_{i-1}).\]
\end{definition}

\clearpage

\begin{theorem}
For any partition $P$,
\[
L(f,P)\le U(f,P).
\]
\end{theorem}

Geometrically:
\begin{itemize}
\item Lower sums inscribe rectangles below the graph,
\item Upper sums circumscribe rectangles above the graph.
\end{itemize}

\clearpage

% ======================================================
\begin{heading}
Refinement and Monotonicity of Sums
\end{heading}

\begin{definition}[Refinement]
A partition $Q$ is a refinement of $P$ if $P\subseteq Q$.
\end{definition}

\clearpage

\begin{theorem}
If $Q$ refines $P$, then
\[
L(f,P)\le L(f,Q)\le U(f,Q)\le U(f,P).
\]
\end{theorem}

Thus:
\begin{itemize}
\item lower sums increase under refinement,
\item upper sums decrease under refinement.
\end{itemize}

This squeezing behavior is the key ton integrability

\clearpage

% ======================================================
\begin{heading}
Upper and Lower Integrals
\end{heading}

Instead of taking limits of sums, Riemann used infima and suprema.

\begin{definition}[Upper Integral]
\[
\overline{\int_a^b} f
=\inf\{U(f,P):P\text{ partition of }[a,b]\}.
\]
\end{definition}

\clearpage

\begin{definition}[Lower Integral]
\[
\underline{\int_a^b} f
=\sup\{L(f,P):P\text{ partition of }[a,b]\}.
\]
\end{definition}

\clearpage

\begin{theorem}
For every bounded function,
\[
\underline{\int_a^b} f\le \overline{\int_a^b} f.
\]
\end{theorem}

\clearpage

% ======================================================
\begin{heading}
Riemann Integrability
\end{heading}

\begin{definition}[Riemann Integrable]
A bounded function $f:[a,b]\to\mathbb{R}$ is \textbf{Riemann integrable} if
\[
\underline{\int_a^b} f=\overline{\int_a^b} f.
\]
The common value is denoted
\[
\int_a^b f(x)\,dx.
\]
\end{definition}

\clearpage

\begin{theorem}[Cauchy Criterion]
$f$ is Riemann integrable if and only if for every $\varepsilon>0$
there exists a partition $P$ such that
\[
U(f,P)-L(f,P)<\varepsilon.
\]
\end{theorem}

This theorem is the practical test for integrability.

\clearpage

% ======================================================
\begin{heading}
Continuity Implies Integrability
\end{heading}

\begin{theorem}
If $f$ is continuous on $[a,b]$, then $f$ is Riemann integrable.
\end{theorem}

\clearpage

% ======================================================
\begin{heading}
Functions with Discontinuities
\end{heading}

Not all integrable functions are continuous.

\begin{problem}
Define
\[
f(x)=
\begin{cases}
1, & x\neq 1,\\
0, & x=1.
\end{cases}
\]
Then $f$ is Riemann integrable on $[0,2]$ and
\[
\int_0^2 f(x)\,dx=2.
\]
\end{problem}

A finite number of discontinuities does not affect area.

\clearpage

% ======================================================
\begin{heading}
A Pathological Counterexample
\end{heading}

\begin{definition}[Dirichlet Function]
\[
f(x)=
\begin{cases}
1, & x\in\mathbb{Q},\\
0, & x\notin\mathbb{Q}.
\end{cases}
\]
\end{definition}

\begin{theorem}
The Dirichlet function is not Riemann integrable on $[0,1]$.
\end{theorem}


\clearpage

% ======================================================
\begin{heading}
Lebesgue's Criterion (Advanced Real Analysis)
\end{heading}

\begin{theorem}[Lebesgue Criterion]
A bounded function is Riemann integrable
if and only if its set of discontinuities has measure zero.
\end{theorem}

This theorem precisely describes the power and limitation of Riemann integration.


\EndPage
\end{document}
